\chapter{Sommario}

Il sistema trasfusionale non può produrre ciò che distribuisce: sangue, plasma
e piastrine dipendono per intero dalla disponibilità volontaria dei donatori.
Capire chi dona, con quale frequenza e per quanti anni continua a farlo è
quindi una premessa operativa alla programmazione della raccolta.
Questa tesi affronta la domanda su dati reali:
l'archivio donazionale del \textbf{Dipartimento Interaziendale di Medicina
Trasfusionale (DIMT)} di Trieste e Gorizia, due province del Friuli Venezia
Giulia con popolazione mediamente più anziana del resto del Paese e una rete
associativa particolarmente radicata.

Il lavoro parte da 27 file Excel eterogenei per schema e convenzioni di
denominazione e li riconduce, attraverso una pipeline in Python documentata in
ogni suo stadio, a un unico dataset analitico di \textbf{353\,671 eventi di
donazione} riferiti a \textbf{43\,299 donatori distinti} nel periodo
\textbf{2009--2025}. Su questa base sono condotte quattro analisi:
un'esplorazione descrittiva della popolazione donante e delle sue tendenze
temporali, rapportate per la prima volta alla popolazione residente; uno studio
di fidelizzazione condotto per coorte di ingresso; una proiezione del volume
donazionale al 2031 selezionata con un backtest a origine mobile fra diciassette
modelli candidati; e un modello di regressione per dati di conteggio sulla
frequenza donazionale individuale. I risultati sono infine resi ispezionabili
in una dashboard interattiva in R Shiny.

Il quadro che emerge è quello di un sistema in contrazione lenta e continua. Il
volume annuale cala del \textbf{28{,}1\%} fra il 2009 e il 2025, concentrando
la flessione nella prima metà del periodo per poi assestarsi; i nuovi ingressi
calano di più; e l'intensità individuale --- il numero di donazioni per
donatore attivo --- resta invariata fra $1{,}84$ e $1{,}97$ lungo tutti e
diciassette gli anni. Il sistema, cioè, non si è intensificato per compensare
il minor ricambio: si è ridotto e poi assestato su una base più piccola, anche grazie
al calo dei consumi.

Rapportare le donazioni alla popolazione residente di Trieste e Gorizia separa
ciò che è demografia da ciò che è comportamento, e il risultato non è quello
che ci si aspetterebbe. Il calo della popolazione in età donabile spiega solo
il \textbf{12{,}2\%} della contrazione; quasi tutto il resto è perdita di
partecipazione vera, da $106{,}9$ a $80{,}1$ donazioni ogni mille residenti. E
la perdita non è dove la
si cercherebbe: a parità di popolazione i \textbf{20--24enni sono la fascia con
il tasso di donazione più alto in ogni anno osservato}, mentre le fasce 35--44
perdono fra il 36 e il 41 per cento. Il problema dei giovani non è che non
cominciano --- sono la metà di tutti gli ingressi --- ma che non tornano: il
\textbf{40{,}6\%} non si ripresenta dopo la prima donazione, contro il
$33{,}3\%$ di chi entra più tardi, e la funzione di rischio decrescente
individua una soglia di consolidamento dell'abitudine intorno al terzo--quinto
anno.

Il modello di frequenza mostra che gli effetti dominanti sono di natura
regolamentare più che motivazionale (genere, IRR $0{,}585$; plasmaferesi, IRR
$1{,}77$: intervalli minimi e tetti annui differenziati) e che, controllando
per gli anni di carriera, \textbf{l'età anagrafica non ha alcun effetto}: ciò
che distingue un donatore assiduo è da quanto tempo dona, non quanti anni ha.
Coerentemente, le proiezioni indicano per il sessennio 2026--2031 circa
\textbf{97\,200 donazioni} complessive: un primo anno sensibilmente sotto il
2025, poi una contrazione di poco meno di un punto percentuale l'anno, senza
ricomposizione attesa fra tipologie. Un assottigliamento lento, non una
rottura.

